\section{Forced Alignment}
\label{app:align}

Given an input audio sequence $X = (x_1,...,x_N)$, where $N$ is the number of frames, a CTC alignment model produces an output sequence $Y = (\bf{y^1,...,y^{T}})$ of length $T$, where $\bf{y^t}$ denotes the posterior probability distribution over an alphabet $\Alphabet$ and blank token $\Blank$. 
Let $\CollapseFunc$ denote the collapsing function of CTC which collapses all repeating symbols and which removes all blanks for a given sequence. 
An alignment path $\boldsymbol{\pi} = \pi_1,...,\pi_T$ in CTC for a given target label of length $M$, $L = (l_1,...,l_M)$ where $\pi_t \in \Alphabet \cup \{\Blank\}$ and $l_i \in \Alphabet$ should satisfy $\CollapseFunc(\boldsymbol{\pi}) = L$. 

Among the all the alignment paths which can be collapsed to the given target label $L$, forced alignment computes the best alignment path $\hat{\boldsymbol{\pi}}$ which maximizes the probability under the posterior distribution given by the acoustic model. 

\begin{equation} \label{eqn:force_align}
\hat{\boldsymbol{\pi}} = \argmax_{\boldsymbol{\pi} \in \CollapseFunc^{-1}(L)} P(\boldsymbol{\pi} | X) 
\end{equation}

CTC assumes every output is conditionally independent of the other outputs given the input. We have 
\begin{equation}
    P(\boldsymbol{\pi} | X) = \prod_{t=1}^{T} P(\boldsymbol{\pi_t} | t, X) = \prod_{t=1}^{T} y^t_{\pi_t}
\end{equation}

\autoref{eqn:force_align} can be computed efficiently using dynamic programming and backtracking.  
An efficient version of this algorithm is implemented in flashlight~\citep{kahn2022flashlight} on CPU. 
We implement a parallel version on CUDA, incorporating memory optimizations that offload memory to CPU, allowing it to process long audio files efficiently as shown in Algorithm~\ref{algo:force_align}.


\begin{algorithm} 
\caption{Pseudo code of our CTC Forced Alignment algorithm on GPU}
\renewcommand{\algorithmicrequire}{\textbf{Input:}}
\renewcommand{\algorithmicensure}{\textbf{Output:}}
\algblockdefx{FORALLP}{ENDFAP}[1]%
  {\textbf{for all }#1 \textbf{do in parallel}}%
  {\textbf{end for}}
\begin{algorithmic}[1]
    \Require{Posterior probabilities $\bf{y}$, target label $L$}
    \Ensure{Alignment path $\pi$} \\ 
    $ S \leftarrow 2 \times |L| + 1$  \\ 
    Create GPU matrices $\alpha_{\text{odd}}$, $\alpha_{\text{even}}$ of size $S$ each to store the maximum log probability of aligning the target upto the given node \\
    B $\leftarrow$ 100; \# buffer size for data copy from GPU to CPU \\ 
    Create CPU matrix $\beta$ of size $T \times S$ for saving the label indices from previous time step used to compute $\alpha_{\text{odd}}$/$\alpha_{\text{even}}$ \\ 
    Create GPU matrix $\beta_{\text{buffer}}$ of size $B \times S$, where B is the buffer size, to store the values of $\beta$ temporarily in a buffer before being copied to CPU. 
    % \zni{should it be updated to (min(B, T), S)?} \vin{This is an implementation detail. The algo works with B x S as well.}
    \For {$t=1,\ldots,T$}
        \FORALLP {$l=1,\ldots,S$} 
            \State If $t$ is odd, update $\alpha_{\text{odd}}[l]$ based on $\alpha_{\text{even}}$, $\bf{y^t}$ using dynamic programming and vice-versa
            \State Store the index $l'$ from previous time step used to update $\alpha_{\text{odd}}$/$\alpha_{\text{even}}$ in  $\beta_{\text{buffer}}$
        \ENDFAP
        \If {t \% B == 0 \textbf{or} t == T}
            \State Copy $\beta_{\text{buffer}}$ to CPU matrix $\beta$ asynchronously
        \EndIf 
    \EndFor
    \\Backtrack and compute $\pi$ from $\alpha_{\text{odd}}$, $\alpha_{\text{even}}$ and $\beta$
\end{algorithmic} 
\label{algo:force_align}
\end{algorithm}



\paragraph{Relationship to Other Alignment Generation Approaches.}
For dealing with noisy transcripts, a commonly used alternative to forced alignment is as follows~\citep{povey2011kaldi}:
first segment the audio into shorter segments that can be input to an acoustic model composed with a language model and generate a transcription for each segment.\footnote{\url{https://github.com/mozilla/DSAlign}} 
Then align the generated transcription with the original text to determine the audio segment for each word. 
If the generated transcriptions cannot be well aligned, then these segments are discarded.

The advantage of this approach is that it enables alignment when the audio and the text do not correspond entirely to each other. 
However, the alignments are performed on a segment per segment basis, whereas our forced alignment process performs a global alignment taking the entire sequence into account. 



% \clearpage
\section{\textit{n}-gram Language Models} 
\label{app:train-lms}

We train 5-gram language models on Common Crawl data using KenLM~\citep{heafield-2011-kenlm} for each language in FLEURS.
For languages that do not use spaces to separate words, we train 20-gram character-level language models.
These languages are Mandarin Chinese (cmn), Cantonese Chinese (yue), Japanese (jpn), Thai (tha), Lao (lao), Burmese (mya) and Khmer (khm).
The text is pre-processed following~\autoref{sec:data-preprocess} and we also remove emojis.\footnote{\url{https://stackoverflow.com/a/33417311}}.

For word-level models, we limit the training data to 40GB and select the 250K most-frequent words as the vocabulary. 
For character-level models, we limit the training data to 6GB.
We provide example commands used for training the LMs below:

\shellcmd{\# word-level LM}\\
\shellcmd{> kenlm/build/bin/lmplz  --prune 1 2 3 4 5 -o 5 --limit\_vocab\_file vocab.txt  -S 90\% -T /tmp/ < input.txt > output.arpa }\footnote{We use \texttt{--prune 1 1 2 3 4} if data size < 5GB and additionally use \texttt{--discount\_fallback} if data size < 1GB }  \\

\shellcmd{\# character-level LM} \\
\shellcmd{> kenlm/build/bin/lmplz  --prune 0 0 0 0 0 1 1 1 2 3 -o 20 trie   -S 90\% -T /tmp/ < input.txt > output.arpa} \\

For n-gram models trained on the FLEURS training data transcriptions, we build 15-gram character level language models without any pruning on all languages in FLEURS. 
For the comparison with Whisper, we only use the Common Crawl language models. 

We use the CTC beam-search decoder from the Flashlight~\citep{kahn2022flashlight} library for decoding our models. 
For decoding with word-level LMs, we use the lexicon-based decoder of~\citet{collobert2016wav2letter, pratap2019w2l} and for character-level LMs, we use the lexicon-free beam-search decoder of~\citet{likhomanenko2019convlm}. 
We tune the language model weight and word insertion penalty on the validation set to select the best hyperparameters for decoding the test set.


\newpage
\section{Comparison to Whisper}
\label{app:asr-whisper}

\autoref{tab:asr-whisper-perlang-test} shows a breakdown of the results into individual languages and \autoref{tab:asr-whisper-full} shows results with and without CC LM n-gram language models.

\insertResultsASRWisperPerLang

\insertResultsASRWisperFull



\clearpage
\section{Comparison to USM}
\label{app:finetune-fleurs-results}

\insertResultsASRUSMFull



% \clearpage
\section{Gender Bias Study}
\label{app:gender-bias-study}

\insertResultsGenderBiasFull
